\section{Discussion \& Conclusion}

\begin{frame}{Discussion~: Vers un récupérateur hybride}
  Pour dépasser le goulot d'étranglement de la récupération~:

  \medskip{}

  \pause[]

  \begin{itemize}
    \item<2-> \textbf{Stratégie hybride~:}
      \begin{enumerate}
        \item<3-> Recherche vectorielle pour identifier des n\oe{}uds
          << graines >>.
        \item<4-> Construction guidée par LLM de la logique de traversée du
          graphe.
      \end{enumerate}
    \item<5-> \textbf{Avantage~:} Se concentrer sur le raisonnement relationnel
      plutôt que sur la syntaxe de la requête.
  \end{itemize}
\end{frame}

\begin{frame}{Conclusion et Perspectives}
  \emph{Med-KAG} confirme la faisabilité de l'ancrage d'un système RAG dans un
  graphe de connaissances médicales.

  \medskip{}

  \pause[]

  \begin{itemize}
    \item<2-> \textbf{Constat~:} L'implémentation actuelle est perfectible
      mais identifie clairement les verrous technologiques.
    \item<3-> \textbf{Travaux futurs~:}
      \begin{itemize}
        \item<4-> Développement d'un récupérateur hybride (en cours).
        \item<5-> Tester la pertinence des réponses avec des experts du milieu
          pour évaluer l'utilité réelle en aide à la décision
        \item<6-> Amélioration de l'explicabilité pour les cliniciens.
      \end{itemize}
  \end{itemize}
\end{frame}

\begin{frame}[standout]
  Merci de nous avoir écouté.

  Des questions?
\end{frame}

